% =========================================================
%  TEMPLATE 2 — MINIMALIST MONOCHROME ASSIGNMENT TEMPLATE
%  Clean black/white academic look with a single accent color
% =========================================================
\documentclass[11pt,a4paper]{article}

\usepackage[margin=1.1in]{geometry}
\usepackage[utf8]{inputenc}
\usepackage[T1]{fontenc}
\usepackage[expansion=false]{microtype}
\usepackage{mathpazo} % elegant serif text/math font

\usepackage[dvipsnames,table]{xcolor}
\definecolor{ink}{HTML}{111111}
\definecolor{accent}{HTML}{B5322C}   % single muted red accent
\definecolor{rule}{HTML}{111111}
\definecolor{codebg}{HTML}{FAFAFA}
\definecolor{coderule}{HTML}{DDDDDD}

\usepackage{titlesec}
\titleformat{\section}
  {\normalfont\Large\bfseries\color{ink}}
  {\thesection.}{0.6em}{}
\titleformat{\subsection}
  {\normalfont\large\bfseries\color{ink}}
  {\thesubsection.}{0.6em}{}
\titleformat{\subsubsection}
  {\normalfont\normalsize\itshape\color{accent}}
  {\thesubsubsection.}{0.6em}{}
\titlespacing*{\section}{0pt}{1.6em}{0.7em}

\usepackage{fancyhdr}
\pagestyle{fancy}
\fancyhf{}
\renewcommand{\headrulewidth}{0.4pt}
\renewcommand{\footrulewidth}{0pt}
\fancyhead[L]{\small\itshape \assignmentTitle}
\fancyhead[R]{\small\itshape \courseCode}
\fancyfoot[C]{\small \thepage}
\fancypagestyle{plain}{\fancyhf{}\renewcommand{\headrulewidth}{0pt}\fancyfoot[C]{\small\thepage}}

\usepackage{graphicx}
\usepackage{amsmath,amssymb}
\usepackage{enumitem}
\usepackage{booktabs}
\usepackage{array}
\usepackage[most]{tcolorbox}
\tcbuselibrary{skins,breakable}
\usepackage{hyperref}
\hypersetup{colorlinks=true,linkcolor=ink,citecolor=accent,urlcolor=accent,pdfborder={0 0 0}}
\usepackage{caption}
\captionsetup{font=small,labelfont={bf}}

\usepackage{listings}
\lstdefinestyle{minimalstyle}{
  backgroundcolor=\color{codebg},
  rulecolor=\color{coderule},
  frame=L,
  framesep=8pt,
  framerule=2pt,
  basicstyle=\ttfamily\footnotesize,
  keywordstyle=\color{accent}\bfseries,
  commentstyle=\color{gray}\itshape,
  stringstyle=\color{ink},
  numberstyle=\tiny\color{gray},
  numbers=left,
  numbersep=8pt,
  breaklines=true,
  showstringspaces=false,
  captionpos=b
}
\lstset{style=minimalstyle}

\newtcolorbox{experimentbox}[1]{
  breakable, colback=white, colframe=ink,
  coltitle=ink, fonttitle=\bfseries\itshape,
  title=#1, boxrule=0pt, toprule=1.2pt, bottomrule=0.6pt,
  leftrule=0pt, rightrule=0pt, sharp corners,
  arc=0mm, left=2pt,right=2pt,top=6pt,bottom=6pt
}
\newcounter{expcounter}
\newenvironment{experiment}[2]{%
  \refstepcounter{expcounter}%
  \begin{experimentbox}{Experiment \theexpcounter \; --- \; #1}
  \textbf{Aim:} \textit{#2} \par\vspace{4pt}
}{\end{experimentbox}}

\newcommand{\assignmentTitle}{Assignment Title Goes Here}
\newcommand{\courseCode}{Course Code}
\newcommand{\courseName}{Full Course Name}
\newcommand{\studentName}{Student Name}
\newcommand{\studentID}{Roll / Registration Number}
\newcommand{\deptName}{Department Name}
\newcommand{\instituteName}{Institute / University Name}
\newcommand{\instructorName}{Instructor Name}
\newcommand{\submissionDate}{\today}
\newcommand{\letterspace}[1]{\spaceskip=0.18em #1}

\begin{document}

% =========================================================
%  TITLE PAGE
% =========================================================
\begin{titlepage}
\newgeometry{margin=1.1in}
\thispagestyle{empty}

\begin{center}
\vspace*{1.2cm}
{\small\MakeUppercase{\letterspace{\instituteName}}}\\[0.15cm]
{\small \deptName}\\[2.2cm]

\rule{0.4\linewidth}{0.8pt}\\[0.6cm]
{\Huge\bfseries \assignmentTitle}\\[0.6cm]
\rule{0.4\linewidth}{0.8pt}\\[1cm]

{\large\itshape \courseCode \; --- \; \courseName}

\vspace{3.5cm}

\begin{minipage}{0.55\linewidth}
\renewcommand{\arraystretch}{1.6}
\begin{tabular}{@{}p{3.2cm}l@{}}
\textit{Submitted by} & \textbf{\studentName} \\
\textit{Student ID}   & \studentID \\
\textit{Submitted to} & \instructorName \\
\textit{Date}         & \submissionDate \\
\end{tabular}
\end{minipage}

\vfill
{\color{accent}\rule{0.15\linewidth}{2pt}}
\end{center}
\end{titlepage}
\restoregeometry

% =========================================================
\tableofcontents
\newpage

\section{Introduction}
Write a brief introduction to the assignment or lab manual here: objective, background,
and scope of the work.

\subsection{Objective}
\begin{itemize}[leftmargin=1.4em]
  \item Objective one.
  \item Objective two.
  \item Objective three.
\end{itemize}

\section{Theory / Background}
Provide the theoretical background, formulas, or concepts relevant to the assignment.

\begin{equation}
  E = mc^2
\end{equation}

\section{Lab Experiments}

Duplicate the \texttt{experiment} block below for every experiment; numbering
increments automatically.

\begin{experiment}{Linear Search Implementation}{To implement and analyze the linear search algorithm.}

\textbf{Procedure:}
\begin{enumerate}[leftmargin=1.4em]
  \item Take an array and a target element as input.
  \item Traverse the array sequentially comparing each element to the target.
  \item Return the index if found, else return $-1$.
\end{enumerate}

\textbf{Source Code:}
\begin{lstlisting}[language=Python, caption={Linear Search in Python}]
def linear_search(arr, target):
    for i, value in enumerate(arr):
        if value == target:
            return i
    return -1

if __name__ == "__main__":
    data = [4, 2, 9, 7, 5]
    result = linear_search(data, 7)
    print(f"Element found at index: {result}")
\end{lstlisting}

\textbf{Output / Observation:}\\
\texttt{Element found at index: 3}

\textbf{Conclusion:}\\
The linear search algorithm successfully locates the target element in
$O(n)$ time complexity in the worst case.

\end{experiment}

\vspace{0.6cm}

\begin{experiment}{Sum of Two Numbers (C Program)}{To write a C program that reads two integers and prints their sum.}

\textbf{Source Code:}
\begin{lstlisting}[language=C, caption={Sum of Two Numbers in C}]
#include <stdio.h>

int main(void) {
    int a, b, sum;
    printf("Enter two integers: ");
    scanf("%d %d", &a, &b);
    sum = a + b;
    printf("Sum = %d\n", sum);
    return 0;
}
\end{lstlisting}

\textbf{Output / Observation:}\\
\texttt{Enter two integers: 5 7}\\
\texttt{Sum = 12}

\end{experiment}

\section{Results and Analysis}
\begin{table}[h]
\centering
\caption{Sample Result Table}
\begin{tabular}{@{}lcc@{}}
\toprule
\textbf{Test Case} & \textbf{Input} & \textbf{Output} \\
\midrule
1 & [4,2,9,7,5], target=7 & Index 3 \\
2 & [1,2,3], target=5     & Not Found \\
\bottomrule
\end{tabular}
\end{table}

\section{Conclusion}
Summarize what was learned or achieved through this assignment / lab exercise.

\section*{References}
\addcontentsline{toc}{section}{References}
\begin{enumerate}[leftmargin=1.4em]
  \item Author, \textit{Title of Book/Resource}, Publisher, Year.
  \item Online resource, \url{https://example.com}, accessed \today.
\end{enumerate}

\end{document}
